%==============================================================================
% Section 15 -- Conclusion
%==============================================================================
\section{Conclusions and Future Directions}
\label{sec:conclusion}

En este estudio comparativo, hemos analizado los fundamentos matemáticos y la integración en sistemas SLAM de 
3D Gaussian Splatting (3DGS)~\cite{kerbl20233d} frente a Neural Radiance Fields (NeRF)~\cite{mildenhall2020nerf}. 
La transición de representaciones implícitas y continuas (NeRF) a representaciones explícitas y discretas (3DGS) 
marca un cambio de paradigma fundamental en los sistemas de percepción robótica, 
trasladando el principal cuello de botella del hardware desde la capacidad de cómputo hacia el consumo de memoria.

Por un lado, 3DGS ha demostrado ser abrumadoramente superior en términos de latencia y velocidad computacional. 
Al reemplazar el costoso proceso de \textit{ray marching} volumétrico de NeRF por una rasterización de primitivas 
explícitas acelerada por GPU, 3DGS logra velocidades de seguimiento y mapeo de 10 a más de 30 
FPS~\cite{keetha2024splatam}, junto con un renderizado fotorrealista en tiempo real que supera los 90 FPS. 
Además, su representación geométrica explícita preserva mejor los detalles de alta frecuencia en las métricas 
de evaluación de mapas (como PSNR y SSIM). Estas características hacen de 3DGS la tecnología idónea para 
aplicaciones interactivas como la realidad aumentada y virtual (AR/VR).

Sin embargo, estas ventajas conllevan compromisos críticos. La naturaleza explícita de 3DGS sufre de un severo 
consumo de memoria de video (VRAM) y espacio de almacenamiento persistente. Mientras que NeRF puede comprimir la
 información de una escena entera en menos de 1 GB de VRAM gracias a los pesos de su red neuronal, 
 3DGS requiere instanciar millones de elipsoides, pudiendo sobrepasar rápidamente los 24 GB en mapeos de 
 trayectorias a gran escala. Adicionalmente, la dependencia de 3DGS de un entorno estrictamente estático hace 
 que la presencia de objetos dinámicos genere artefactos visuales, requiriendo el uso de complejas extensiones 
 temporales como 4DGS que hoy en día resultan computacionalmente prohibitivas.

En resumen, 3DGS-SLAM unifica la reconstrucción fotorrealista de alta fidelidad y el seguimiento de cámaras en 
tiempo real, superando las barreras de latencia de las arquitecturas NeRF. El éxito y la adopción futura de 3DGS 
en plataformas con recursos limitados (como robots móviles o sistemas integrados) dependerán de la investigación 
activa en técnicas agresivas de compresión, como la cuantización vectorial, la poda de atributos y la codificación
de entropía, así como de un modelado eficiente para escenas dinámicas.